\chapter{Approche avancée – du clustering au scoring continu}

\section*{Introduction}
Suite aux limites découvertes avec l'approche par partitionnement au Chapitre 3, ce chapitre documente notre second cycle expérimental qui reformule le problème pour optimiser la mesure de la performance ESG. Nous détaillons ici l'échec des méthodes d'agrégation simplistes, l'analyse du biais géographique comme obstacle à la segmentation, et la nécessité de construire un score maison capable de synthétiser le désaccord entre agences.

\section{Redéfinition de la compréhension métier : du clustering au scoring}

La première étape du second cycle a consisté à redéfinir l'objectif central du projet et ses critères de succès pour se concentrer sur une approche de scoring continu.

\subsection{Pivot de l'objectif : de la « catégorie » à la « performance relative »}

La performance ESG est un spectre continu, non une collection de catégories. Ce constat est étayé par la littérature sur le « rating disagreement » \citep{bissoondoyal2024}, qui démontre que les agences (MSCI, Sustainalytics, Asset4) divergent radicalement dans leurs notations. Puisque les experts ne s'entendent sur aucune frontière stable, notre objectif devient d'entraîner un modèle de scoring capable de situer chaque entreprise sur un continuum de 0 à 100.

\subsection{Le problème du biais géographique}

Avant de construire ce score, nous devons affronter un problème structurel : le pays d'origine. La carte choroplèthe ci-dessous, représentant l'intensité carbone moyenne (Scope 1 / CA) par pays pour nos 13~518 entreprises, illustre cette fracture.

\begin{figure}[H]
\centering
\includegraphics[width=0.8\textwidth]{images/intensite_carbone_pays.png}
\caption{Intensité carbonne}
\label{fig:Intensité carbonne}
\end{figure}

L'analyse de cette carte montre que les disparités ne reflètent pas forcément une mauvaise gestion individuelle, mais des contextes nationaux. Ignorer cette structure reviendrait à punir une entreprise pour son environnement macroéconomique. Le pays ne doit donc pas être une simple étiquette, mais une variable de variance à traiter.

\section{Diagnostic des échecs du scoring traditionnel}

Pour répondre à ce défi, nous avons mis à l'épreuve deux méthodes classiques afin de vérifier si une agrégation simple suffit à absorber cette complexité géographique.

\subsection{Tentative 1 : la moyenne simple}

La méthode la plus élémentaire consiste à calculer la moyenne arithmétique des variables d'intensité standardisées. Ce score est rapide à produire et facile à interpréter.

\begin{equation}
    Score_i = \frac{1}{n} \sum_{j=1}^{n} x_{i,j}
    \label{eq:score_moyenne_simple}
\end{equation}
\noindent\textit{Limite fondamentale : cette approche suppose que toutes les variables sont indépendantes, ignorant les fortes corrélations démontrées au Chapitre~2.}

\subsection{Tentative 2 : la pondération arbitraire}

Une deuxième approche consiste à attribuer des poids prédéfinis aux trois piliers, afin de refléter des priorités stratégiques. Mathématiquement, le score d'une entreprise $i$ est une somme pondérée des moyennes de chaque pilier :

\begin{equation}
    Score_i = w_E \cdot \bar{E}_i + w_S \cdot \bar{S}_i + w_G \cdot \bar{G}_i
    \label{eq:score_pondere}
\end{equation}

\noindent\textit{Limite fondamentale : elle est hautement subjective. Le score final dépend uniquement du choix personnel de l'analyste plutôt que de la réalité mathématique des chiffres. Cela ne fait que reproduire le désaccord entre agences au lieu de le résoudre par les données.}

\subsection{Échec commun : écrasement de la variance}

Le graphique de densité (KDE) ci-après montre que ces deux méthodes présentent un défaut commun : elles écrasent la variance des données. En additionnant des variables sans tenir compte de leurs corrélations, elles produisent des distributions resserrées où il devient difficile de distinguer les entreprises réellement performantes des moins performantes.

\begin{figure}[H]
\centering
\includegraphics[width=0.8\textwidth]{images/Comparaison des scores simplistes.png}
\caption{Comparaison des scores simplistes}
\label{fig:comparaison des scores simplistes}
\end{figure}
\section{Solution : L'Analyse en Composantes Principales (PCA) Hiérarchique}

Face aux limites des méthodes classiques de scoring ESG – subjectivité des pondérations, biais géographique – nous avons développé une approche fondée sur une PCA hiérarchique à deux niveaux.

\subsection{Nettoyage des colinéarités : une étape préalable indispensable}

Avant d’appliquer la PCA, nous devons supprimer les redondances entre variables. Selon O’Brien (2007), une colinéarité excessive gonfle artificiellement la variance des poids, rendant le score instable. Pour pallier ce biais, nous avons appliqué une stratégie de réduction de dimension en deux étapes :

Filtrage par corrélation ($r > 0,8$) : Élimine les doublons directs pour ne pas surestimer un seul indicateur.

VIF (Variance Inflation Factor) itératif : mesure à quel point une variable est expliquée par les autres variables du même groupe. Nous avons fixé un seuil de VIF < 10, limite critique identifiée par Hair et al. (2010) pour garantir l'indépendance statistique.

Cette méthode a permis de condenser le pilier E (12 variables à 4) et le pilier S (4 variables à 2), tandis que le pilier G est resté intact (4 variables) .Cette sélection garantit que chaque variable apporte une information unique, ce qui stabilise les étapes suivantes.

\subsection{PCA de niveau 1 : synthétiser chaque pilier}

L’enjeu est de transformer les indicateurs techniques en trois scores simplifiés, un pour chaque pilier (E, S et G). Pour ce faire, nous utilisons une PCA de Niveau 1.

Contrairement aux agences de notation qui fixent des coefficients de manière arbitraire, la PCA laisse les données "parler". Selon Jolliffe (2002) ( doit insérer la référence : Jolliffe, I. T. (2002). Principal Component Analysis. Springer Series in Statistics. 2nd Edition.) , cette méthode permet d'extraire les axes de variance maximale. En d'autres termes, le poids de chaque variable dans le score est dicté par sa capacité réelle à différencier les entreprises, éliminant ainsi toute subjectivité humaine.

Le tableau ci-dessous retrace le passage des données brutes vers les scores synthétiques.

\begin{table}[H]   % H majuscule = ici exactement
\centering
\renewcommand{\arraystretch}{1.5}
\small
\resizebox{\textwidth}{!}{%
\begin{tabular}{|p{2.6cm}|c|c|c|c|p{2.9cm}|}
\hline
\textbf{Pilier} & 
\thead{\textbf{Variables}\\\textbf{Initiales}} & 
\thead{\textbf{Après}\\\textbf{Nettoyage}\\\textbf{(VIF/Corr)}} & 
\thead{\textbf{Synthèse}\\\textbf{PCA}\\\textbf{(Composantes)}} & 
\thead{\textbf{Information}\\\textbf{Conservée}} & 
\centering\arraybackslash\thead{\textbf{Vecteur}\\\textbf{Principal}\\\textbf{(Loading)}} \\
\hline
Environnement (E) & 12 & 4 & 3 & 92,57\% & Émissions CO2 (Scope 1 \& 3) \\
\hline
Social (S) & 4 & 2 & 2 & 100,00\% & Capital Humain / Formation \\
\hline
Gouvernance (G) & 4 & 4 & 1 & 83,78\% & Structure de Revenus \& Cap. Boursière \\
\hline
\end{tabular}%
}
\normalsize
\end{table}
\textbf{Indépendance des scores obtenus} 
\begin{figure}[H]
\centering
\includegraphics[width=0.8\textwidth]{images/corr_pillars_v3.png}
\caption{Corrélations entre les scores des piliers ESG}
\label{fig:Corrélations entre les scores des piliers ESG}
\end{figure}
La robustesse de cette étape est validée par la faible corrélation entre les scores obtenus. Cela prouve que chaque pilier mesure une dimension unique du risque ESG, garantissant un score final équilibré.
\subsection{PCA de niveau 2 : la pondération naturelle}

L’étape finale consiste à agréger les scores E, S et G en un Score ESG Global. Plutôt que d'attribuer des poids arbitraires, nous appliquons une PCA de second niveau sur des données préalablement standardisées, pour laisser émerger la structure réelle des données.

Le graphique ci-dessous présente les poids extraits de la première composante principale (PC1), sélectionnée selon le critère de Kaiser (valeur propre > 1) :
\begin{figure}[H]
\centering
\includegraphics[width=0.8\textwidth]{images/pillars_weights_v3.png}
\caption{Contribution des piliers au score ESG global}
\label{fig:Contribution des piliers au score ESG global}
\end{figure}
L’environnement domine, suivi du social, puis de la gouvernance. Cette hiérarchie est issue des corrélations observées dans nos données. Aucune intervention humaine n’a biaisé ces coefficients.

\subsubsection{Normalisation par rangs : une distribution uniforme pour le machine learning}

Le score brut issu de la PCA de niveau 2, suit une loi normale où la majorité des entreprises se concentrent autour de la moyenne. Ce qui réduit le pouvoir discriminant nécessaire aux algorithmes prédictifs. Pour y remédier, nous appliquons une normalisation par rangs via l’outil QuantileTransformer. Cette technique redistribue les scores de façon homogène sur un intervalle de [0, 100].
\begin{figure}[H]
\centering
\includegraphics[width=0.8\textwidth]{images/score_distribution_comparison_v3.png}
\caption{Comparaison de la distribution du score}
\label{fig:Comparaison de la distribution du score}
\end{figure}
Comme l’illustre la comparaison des histogrammes avant et après traitement, nous passons à une distribution uniforme. Avec une moyenne de 50,44 et des quartiles parfaitement équilibrés (25e : 25,4 ; 75e : 75,7).

\subsubsection{Neutralisation géographique : comparer des pairs locaux}

Les entreprises ne sont pas exposées aux mêmes contraintes environnementales selon leur pays. Pour éviter qu’une entreprise ne soit pénalisée simplement à cause de sa localisation, nous centrons-réduisons le score par pays, puis nous appliquons un nouveau QuantileTransformer global. 

\begin{figure}[H]
\centering
\includegraphics[width=0.8\textwidth]{images/score_by_country_before_v3.png}
\caption{Distribution du score ESG par pays avant neutralisation}
\label{fig:Distribution du score ESG par pays avant neutralisation}
\end{figure}
\begin{figure}[H]
\centering
\includegraphics[width=0.8\textwidth]{images/score_by_country_after_v3.png}
\caption{Distribution du score ESG par pays aprés neutralisation}
\label{fig:Distribution du score ESG par pays aprés neutralisation}
\end{figure}



Le premier graphique montre que les scores varient naturellement selon le pays d'origine, créant un biais géographique injuste. Le second graphique prouve que la méthode de centrage-réduction, préconisée par Zou et al. (2023), élimine ce biais en alignant tous les pays sur une base commune de 50.\cite{zou2025impact}


\section{Robustesse Statistique du Score pour le Machine Learning}

Cette étape finale valide la viabilité du score ESG comme variable cible pour les modèles prédictifs. 

\subsection{Analyse des corrélations}

La corrélation finale confirme la dominance du pilier Environnement dans la constitution de la performance globale. La baisse de la corrélation du pilier Social par rapport aux poids initiaux de la PCA s'explique par la neutralisation géographique : une part significative de la variance de ce pilier était liée aux standards nationaux plutôt qu'aux efforts individuels des entreprises.

\begin{figure}[H]
\centering
\includegraphics[width=0.8\textwidth]{images/final_pillar_contribution.png}
\caption{Contribution finale des piliers au score ESG global}
\label{fig:Contribution finale des piliers au score ESG global}
\end{figure}
\subsection{Validation de la distribution}

L'analyse de la distribution confirme que notre distribution suit une loi uniforme, assurant ainsi une densité de données équivalente sur tout l'intervalle [0, 100].

Skewness (Asymétrie) : -0.0134  (Distribution parfaitement symétrique).  
Kurtosis : -1.2108 (Aplatissement caractéristique d'une distribution uniforme).

\subsection{Pouvoir discriminant : Courbe de Lorenz}

La courbe de Lorenz démontre que le score maximise l’entropie en épousant la diagonale d'égalité parfaite, garantissant ainsi un pouvoir discriminant optimal où chaque entreprise occupe une position unique pour la modélisation future.

\begin{figure}[H]
\centering
\includegraphics[width=0.8\textwidth]{images/validation_lorenz_v3.png}
\caption{Capacité de discrimination}
\label{fig:Capacité de discrimination}
\end{figure}

\subsection*{Conclusion}

En combinant la synthèse par PCA, la neutralisation des biais géographiques et une normalisation par rangs, nous avons transformé des données brutes hétérogènes en un indicateur robuste, équitable et statistiquement optimisé pour la phase de modélisation prédictive.
